% !TEX program = lualatex
% トヨタ関連企業 研究者データベース構築の進捗報告（今井泰介先生宛）
% 組版は LuaLaTeX。詳細と既知の不具合は同ディレクトリの README.md を参照。
\documentclass[a4paper,10pt]{ltjsarticle}

\usepackage{booktabs}
\usepackage[margin=22mm]{geometry}

% ltjsarticle の \paragraph は既定で ■ を見出し記号として出力する。
% この報告書では不要なので、クラスが用意している差し替え点を空にする。
% （titlesec で見出し書式を上書きすると ■ が本文へ落ちるため使わない）
\makeatletter
\renewcommand{\jsParagraphMark}{}
\makeatother

\title{トヨタ関連企業 研究者データベースの構築状況}
\author{伊倉涼介}
\date{2026年8月17日}

\begin{document}
\maketitle


\section{本報告について}

「誰が、どの論文を、誰と書いたのか」を網羅したデータベースは、
本共同研究Science of Science の分析課題の前提となる。
本報告ではその構築状況と分かった事項をまとめて報告する。

\section{データベースの構築状況}

トヨタの三組織の公式サイトから論文一覧919件を自動取得し、
うち852件（92.7\%）を OpenAlexに
対応づけた（内訳は表\ref{tab:corpus}。AISIN IMRA は日本法人と米国法人で別サイト）。

\begin{table}[htbp]
\centering
\caption{公式サイトから収集した論文と OpenAlex への対応づけ状況}
\label{tab:corpus}
\begin{tabular}{lrrr}
\toprule
収集元 & 収集件数 & 対応づけ済 & 対応率 \\
\midrule
豊田中央研究所 & 355 & 352 & 99.2\% \\
AISIN IMRA（米国） & 313 & 292 & 93.3\% \\
トヨタ自動車未来創成センター & 154 & 131 & 85.1\% \\
AISIN IMRA（日本） & 97 & 77 & 79.4\% \\
\midrule
合計 & 919 & 852 & 92.7\% \\
\bottomrule
\end{tabular}
\end{table}

DOIがある論文は完全一致で、ない論文はタイトルの類似度で対応づけた。
文字列一致だけでは誤同定が起こりうるため、判断のつかない10件を目視で確認したところ、
9件は同一論文であったが、1件は掲載誌も出版年も異なる別論文であったため取り消した（表の852件は取消後）。
対応づかなかった67件のうち66件は DOI を持たず、OpenAlex 未収録と見られる。

公式サイトに一覧のない AISIN IMRA 欧州法人は、
Web of Science（大阪大学が契約する商用の文献データベース）の機関名検索で25件を得た。
同じ検索は米国法人とアイシン精機で833件を返し、
うち683件は公式サイトに掲載がなかった。機関名検索は公式サイトの掲載漏れを補う手段になる。

被引用24,786件と参考文献27,052件も取得し、指標計算に使える。
著者を研究者単位にまとめ直すと1,249名で、うち1,134名（90.8\%）の所属が判明し、
574名はトヨタ関連企業への所属を確認できた。

\section{使用したデータベースとその使い分け}

構築には四つの外部データベースを用いており、その分担は一つの判断から導かれている。
\textbf{著者を名前で検索するより、論文から引く方が安全である}という判断である。
著者名で検索すると同姓同名が必ず混ざり、名前だけを手がかりにする限りこの曖昧さは解消できない。
一方、論文は DOI で一意に定まる。
どの論文かが確定していれば、そこから著者名や所属を取り出しても人物の取り違えは起きない。
すなわち「この人物は誰か」という難しい問いを、
「この論文の著者は誰か」という確実に答えられる問いに置き換えている。

\subsection{OpenAlex（無料）　論文の同定と指標計算の土台}

OpenAlex からは、論文の同定と、指標計算の土台になるデータ（被引用、参考文献、分野分類）を取る。
一方で\textbf{著者に振られた識別番号は使わない}。
同姓同名の別人が一つの番号にまとめられているためで、その実例は第4.3節に述べる。
無料枠は1日1,000回の取得に限られるが、1回で50件をまとめて取れるため実質1日5万件にあたり、
今回の規模には足りる。

\subsection{Crossref（無料）　著者名と ORCID の補完}

Crossref は著者名と ORCID を補う。
\textbf{AISIN IMRA 米国法人の313件は公式サイトに著者名の記載が一切なかった}ため、ここで埋めた。
ORCID は383名分を得た。第4.3節に述べるとおり、著者に紐づく識別子のうち信頼できるのはこれだけである。
弱点は所属情報にある。Crossref が所属を持っていた論文は34.5\%にとどまる。
出版社が登録した場合しか入らないためである。この不足を埋めるのが Scopus である。

\subsection{Scopus（大阪大学の契約）　論文を経由した所属の取得}

Scopus からは著者の所属を取る。
ここには方針を変えた経緯があり、それ自体が本節冒頭の判断の裏づけになっている。
当初は著者を名前で検索し、その著者に紐づく所属を直接得るつもりであったが、二つの理由で断念した。
一つは、ORCID による検索が試した6名すべてで0件だったことである。
Scopus 側に ORCID がほとんど登録されていない。
もう一つは、名前での検索が同姓同名を排除できないことである。
「Suzuki」という姓では\textbf{45,844件}が該当し、名前だけでは一人に絞り込めない。

そこで、論文から所属を引く方式に切り替えた。どの論文かは確定しているため、取り違えは起きない。
結果として所属が分かる論文の割合は\textbf{34.5\%から96.1\%}へ上がり、
Scopus だけが所属を持っていた論文も\textbf{394件}あった。
第2節の1,134名の所属が判明したという結果は、この方式によるものである。

\subsection{Web of Science（大阪大学の契約）　機関単位の探索}

Web of Science は他の三つと違い、機関を単位に引く。
\textbf{返ってくるデータに所属の欄がない}ため、
個々の論文について著者の所属を確かめる用途には使えない。
代わりに機関名での絞り込みが効くので、機関単位で論文を洗い出す用途に限定している。
第2節で述べた欧州法人の回収と、公式サイトの掲載漏れの補完がこれにあたる。
予備的な検証では著者名と機関名の組み合わせを所属の間接証拠に使えるかを試したが、
\textbf{Scopus が見つけられなかった所属を1件も新規発見しなかった}ため、この用途からは外した。

なお費用について、OpenAlex と Crossref は無料の範囲に収まり、
Scopus と Web of Science は大阪大学の契約の範囲にある。追加の費用は発生しない。

\section{構築を通じて分かった事項}

\subsection{収録している年代がサイトによって大きく異なる}

出版年が2018年以降の割合は、未来創成センター98.5\%、豊田中央研究所99.4\%に対し、
AISIN IMRA 米国法人17.8\%（中央値2006年、292件中62件は2000年より前）、
日本法人51.9\%、全体67.0\%である。
期間の取り方によって、どの組織のデータがどれだけ残るかが変わる。

\subsection{受領した著者リストは裏付けの取れる割合が高くない}

受領した著者リスト（Sota Hayashi 氏より9,429名分、OpenAlex 由来）の100名を
Scopus（商用の文献データベース）と Web of Science で照合したところ、
裏付けの取れた割合は異常なキャリア長の群30.0\%、
「トヨタ自動車のスイス法人」と登録された組織の群48.6\%、対照群62.9\%であった。
スイス法人の群の11名（31.4\%）は実際には東京大学や三菱総合研究所などに所属し、
抽出元によっては母集団に使えない。

\subsection{著者の識別番号が信頼できないため、名寄せの誤り率を実測した}

受領リストの精度が伸びない原因もここにある。
同一人物が複数の表記で現れるため、これを一人に統合する処理（名寄せ）が最も難しい。
受領リストと同じ OpenAlex の著者識別番号は\textbf{信頼できない}。
同姓同名の別人が一つの番号にまとめられ、「S. Nakamura」は研究歴が74年に及ぶ。
Web of Science の研究者識別番号も同様で、2022年発行の番号に1970年代の論文が対応づく。
研究歴が異常に長い著者はこの集約が原因のため（該当群の76.7\%が同姓同名の多い名前、他は14〜17\%）、
除外せず複数人物として分離する。

一方 ORCID（研究者が自ら登録する国際的な識別番号）は良好で、
383名のうち368名が一つの氏名表記のみに対応した。
そこで ORCID を主たる手がかりとし、持たない著者は姓と名の頭文字で名寄せし、
ORCID 保有者を正解に誤り率を実測した。
別人を同一人物とする誤りが4.3\%、同一人物を分ける誤りが1.0\%である。
ただし ORCID 保有者はフルネームが記録された層に偏り、頭文字のみの著者では誤り率がこれより悪い可能性がある。
頭文字しかなく ORCID もない154名は目視確認に回す。

\subsection{組織を横断して研究する著者を21名検出した}

たとえば Tomohiko Jimbo 氏は未来創成センターと豊田中央研究所の双方に、
Norihiro Mizoshita 氏は AISIN IMRA（日本）と豊田中央研究所の双方に論文がある。
これら21名（うち12名は ORCID で裏付け）は、
先生が挙げられた「どのようなチーム構造が良い研究成果につながるか」の分析対象になり得る。
ただし検出までの段階であり、設計は以下の判断に依る。

\section{判断を仰ぎたい事項}

\paragraph{一つめ　分析の対象期間}
対象期間を、未来創成センターと豊田中央研究所が揃う2018年以降に限るか全期間とするか決めたい。
前者は AISIN IMRA 米国法人の8割を失い、後者も Web of Science で
補える2000年より前が121件と薄い。
ただしご指示の hot streak のようなキャリア分析に10年未満では足りず、
キャリアを追うのであれば全期間を提案する。

\paragraph{二つめ　主として用いるデータ}
主たるデータを、裏付けの取れる割合が30.0〜62.9\%にとどまる受領リストとするか、
独自構築のデータベースとするか決めたい。
\textbf{後者を提案する。}トヨタ自身の公表情報に基づき、誤り率も実測できているためである。
ただし収録範囲は三組織に限られ、年代の偏りは一つめの判断に依存する。

\paragraph{三つめ　対象とする組織の範囲}
三組織をどこまで広げるかも決めたい。トヨタ関連とみなしうる組織の論文数は、
Web of Science で豊田工業大学6,295件、デンソー1,457件、
Toyota Research Institute 658件、JTEKT 279件、Toyota Motor Europe 278件、豊田自動織機127件などである。
豊田工業大学は教育研究機関であり、企業内研究を扱う本研究の趣旨から外れるため除外を提案する。
残りは第2節と同じ機関名検索で回収できる。
まず現在の三組織で分析を始め、必要に応じて広げることを提案する。

\section{今後の計画}

所属情報は同一組織が住所付きの正式名称と略称で別々に記録されているため、
まず表記を揃えて組織ごとに集計できる状態にし、
独自データベースを基準に受領リストの精度を数値で示す。
一つめの判断で対象期間を絞り込み、三つめの判断に応じて機関名検索の範囲を広げる。
三点が定まってから指標計算の設計に入る。
第4.2節の検証の詳細は別途レポートに、データと処理手順は非公開リポジトリにまとめている。

\end{document}
